\SourcePage{288}{293}
\section{Kinematic invariants}

Let us consider some kinematic relations for scattering processes, in which both in the initial and in the final states there are just two particles each. We have in mind relations which are consequences of the general conservation laws alone and are therefore valid irrespective of the nature of the particles and of the laws of their interaction.

\SourcePage{289}{294}
Let us write the law of conservation of 4-momentum in a general form, not specifying which of the momenta refer to the initial particles and which to the final ones:
\begin{equation}
q_1 + q_2 + q_3 + q_4 = 0.
\tag{66.1}
\end{equation}
Here $\pm q_a$ are the 4-vectors of momenta, two of them corresponding to the incoming particles, and two to the scattered ones; for the latter the momenta are $-q_a$.  In other words, two of the $q_a$ have temporal component $q^0_a > 0$, and two have $q^0_a < 0$.

Alongside the law of conservation of 4-momentum, the law of conservation of charge should also be obeyed. By charge one should understand here not only the electric charge, but also other conserved quantities, having different sign for particles and antiparticles.

From the given species of particles participating in the process, the squares of the 4-vectors $q_a$ are the given squared masses of the particles ($q^2_a = m^2_a$). Depending on the values of the time components $q^0_a$ run through, and on the values of the charges, we obtain three different reactions. Let us write these three processes as:
\begin{equation}
\begin{aligned}
&\text{I.}\quad 1 + 2 \to 3 + 4,\\[3pt]
&\text{II.}\quad 1 + \bar{3} \to \bar{2} + 4,\\[3pt]
&\text{III.}\quad 1 + \bar{4} \to \bar{2} + 3.
\end{aligned}
\tag{66.2}
\end{equation}
Here the digit indicates the number of the particle, and a bar over it indicates the antiparticle. The transition from one of the reactions to another, i.e.\ the transfer of a particle from one side of the formula to the other, corresponds to a change of sign of the temporal component $q^0_a$, and also of the sign of the charge, i.e.\ the replacement of the particle by the antiparticle. (Alongside the processes~(66.2), the inverse reactions are possible, of course.)

The three processes~(66.2) are referred to as three \textit{crossed} (or \textit{cross-channels}) of one (generalised) reaction.

Let us give several examples. If particles~1 and 3 are electrons, and 2 and 4 are photons, then channel~I represents scattering of a photon by an electron; by virtue of the true neutrality of the photon channel~III represents the same as~I. Channel~II is the conversion of an electron-positron pair into two photons. If all four particles are electrons, then channel~I is electron-electron scattering, and channels~II and~III are electron-positron scattering. If particles~1 and 3 are electrons, and 2 and 4 are muons, then channel~I is $e\mu$ scattering, channel~III is $e\bar{\mu}$ scattering, and channel~II is the conversion of the pair $e\bar{e}$ into a pair $\mu\bar{\mu}$.

In the study of scattering processes an especially important role is played by invariant quantities, which can be composed from the 4-momenta. They are functions of the invariant amplitudes of scattering (see~\S\,70).

\SourcePage{290}{295}
From the four 4-momenta one can compose two independent invariant quantities. Indeed, by virtue of~(66.1) only three of the four 4-vectors $q_a$ are independent; let these be $q_1$, $q_2$, $q_3$. From them one can compose six invariants: three squares $q^2_1$, $q^2_2$, $q^2_3$ and three products $q_1 q_2$, $q_1 q_3$, $q_2 q_3$. But the first three are the given squared masses, and the latter three are connected by one relation, following from the equality
\[
(q_1 + q_2 + q_3)^2 = q^2_4 = m^2_4.
\]

For the sake of greater symmetry it is convenient, however, to consider not two but three invariants, which we choose as follows\,\footnote{In the general case, when the reaction involves $n \geqslant 4$ particles, the number of functionally independent invariant variables is $3n - 10$. Indeed, there are in all $4n$ variables---components of $n$ 4-momenta $q_a$. Among them there are $n$ functional relations $q^2_a = m^2_a$, and four more given by the law of conservation $\sum q_a = 0$. Arbitrary values can also be given to six parameters, determining the general Lorentz transformation (a general four-dimensional rotation). Therefore the number of independent invariant variables is: $4n - n - 4 - 6 = 3n - 10$.}:
\begin{equation}
\begin{aligned}
s &= (q_1 + q_2)^2 = (q_3 + q_4)^2,\\[3pt]
t &= (q_1 + q_3)^2 = (q_2 + q_4)^2,\\[3pt]
u &= (q_1 + q_4)^2 = (q_2 + q_3)^2.
\end{aligned}
\tag{66.3}
\end{equation}
They are connected, as is easy to see, by the relation
\begin{equation}
s + t + u = h,
\tag{66.4}
\end{equation}
where
\begin{equation}
h = m^2_1 + m^2_2 + m^2_3 + m^2_4.
\tag{66.5}
\end{equation}

In channel~I the invariant $s$ has a simple physical meaning. It is the square of the total energy of the colliding particles (1 and~2) in the system of their centre of inertia (with $\mathbf{p}_1 + \mathbf{p}_2 = 0$: $s = (\varepsilon_1 + \varepsilon_2)^2$). In channel~II an analogous role is played by the invariant~$t$, and in channel~III the invariant~$u$. In connection with this, channels~I, II, III are often called the $s$-, $t$-, and $u$-channels.

It is not difficult to express the invariants $s$, $t$, $u$ through the energies and momenta of the colliding particles in each of the channels. Let us consider the $s$-channel. In the centre of inertia system of particles 1 and~2 the temporal and spatial components of the 4-vectors $q_a$ are given as follows:
\begin{equation}
\begin{aligned}
q_1 &= p_1 = (\varepsilon_1, \mathbf{p}_s),\qquad &q_2 &= p_2 = (\varepsilon_2, -\mathbf{p}_s),\\[3pt]
q_3 &= -p_3 = (-\varepsilon_3, -\mathbf{p}'_s)\qquad &q_4 &= -p_4 = (-\varepsilon_4, \mathbf{p}'_s)
\end{aligned}
\tag{66.6}
\end{equation}

\SourcePage{291}{296}
(the index $s$ on $\mathbf{p}_s$, $\mathbf{p}'_s$ indicates that these momenta refer to the reaction in the $s$-channel). Then
\begin{equation}
s = \varepsilon^2_s,\qquad \varepsilon_s = \varepsilon_1 + \varepsilon_2 = \varepsilon_3 + \varepsilon_4;
\tag{66.7}
\end{equation}
\begin{equation}
4\mathbf{p}^2_s = [s - (m_1 + m_2)^2][s - (m_1 - m_2)^2],
\tag{66.8}
\end{equation}
\begin{equation}
4\mathbf{p}'^2_s = [s - (m_3 + m_4)^2][s - (m_3 - m_4)^2];
\end{equation}
\begin{equation}
2t = h - s + 4\mathbf{p}_s\mathbf{p}'_s - \frac{1}{s}(m^2_1 - m^2_2)(m^2_3 - m^2_4),
\tag{66.9}
\end{equation}
\[
2u = h - s - 4\mathbf{p}_s\mathbf{p}'_s + \frac{1}{s}(m^2_1 - m^2_2)(m^2_3 - m^2_4).
\]

In the case of elastic scattering ($m_1 = m_3$, $m_2 = m_4$) we have $|\mathbf{p}_s| = |\mathbf{p}'_s|$, so that $\varepsilon_1 = \varepsilon_3$, $\varepsilon_2 = \varepsilon_4$. Instead of~(66.9) the simpler formulae then hold:
\begin{equation}
\begin{aligned}
t &= -(\mathbf{p}_s - \mathbf{p}'_s)^2 = -2\mathbf{p}^2_s(1 - \cos\theta_s),\\[3pt]
u &= -2\mathbf{p}^2_s(1 + \cos\theta_s) + (\varepsilon_1 - \varepsilon_2)^2,
\end{aligned}
\tag{66.10}
\end{equation}
where $\theta_s$ is the angle between $\mathbf{p}_s$ and $\mathbf{p}'_s$. Let us note that the invariant $-t$ represents the square of the (three-dimensional) momentum transferred in the collision.

Analogous formulae for other channels are obtained by a simple change of notation. For the transition to the $t$-channel one should make the substitution $s \leftrightarrow t$, $2 \leftrightarrow 3$; for the transition to the $u$-channel the substitution $s \leftrightarrow u$, $2 \leftrightarrow 4$.
